\chapter{ImProver Experiments}\label{chap:improver-exp}

We test ImProver on rewriting real-world undergraduate theorems, competition problems, and research-level mathematics and compare its results to those of the base GPT-4o and GPT-4o-mini models. We examine the optimization capabilities of ImProver for the length and declarative metrics, studying the effectiveness in maintaining the correctness of the tactic proof while making it more concise as well as making it more declarative in style and structure.

\subsection{Setup}
\label{sec:improver_setup}

Our experimentation is split into three distinct stages. We first perform ablation testing on the ImProver model parameters (\S\ref{ssec:improver_sampling}) to ensure that ImProver's parameter specification is the optimal one with respect to correctness and metric optimization score. We then evaluate this optimal parameter combination on datasets of varying complexity and analyze the performance and results thereof. Lastly, we note the performance of ImProver in NTP applications in comparison to the base GPT-4o and GPT-4o-mini models.

\subsubsection{Datasets}

We evaluate ImProver on subsets of three datasets of varying difficulty:

\begin{itemize}
    \item \textbf{Mathematics in Lean (MIL)} \cite{MIL} --- 72 theorems drawn from chapters covering set theory, number theory, group theory, topology, calculus, and integration. These are undergraduate-level results stated and proved in a clear, declarative style but often verbose.
    \item \textbf{Compfiles} \cite{compfiles} --- 26 theorems drawn from IMO/USAMO competition problems from 2016--2024. These represent a higher level of mathematical difficulty.
    \item \textbf{Mathlib} \cite{The_mathlib_Community_2020} --- 43 theorems from research-level algebraic topology (specifically, the FundamentalGroupoid component). These are the most difficult and represent active research-level mathematics.
\end{itemize}

Additionally, ablations are performed on a subset of MIL, and theorem proving is benchmarked on various subsets of MIL as well as the MiniF2F \cite{zheng2022miniff} dataset.

\subsubsection{Models and Baseline}

Our base generator uses GPT-4o \cite{openai2024gpt4technicalreport} (\texttt{gpt-4o-2024-08-06}). Since no prior methods currently exist for automated proof optimization, we consider a prompted GPT-4o without the improvements described in Chapter~\ref{chap:improver} as our baseline. Additionally, the baseline and ImProver both receive a prompt containing instructions to optimize for the given metric, with the theorem statement, context, and initial proof. ImProver augments this prompt with the data from the improvements described in Chapter~\ref{chap:improver}.

\subsubsection{Performance Metrics}

Since proof optimization is a new task, we define four performance metrics for measuring aspects of correctness and improvement.

First, we define \textit{improvement} for length as percentage change in length, $\frac{\mu_{\text{len}}(y_0)-\mu_{\text{len}}(y)}{\mu_{\text{len}}(y_0)}\times 100$. For readability, we use the difference, $\mu_{\text{read}}(y)-\mu_{\text{read}}(y_0)$. If no correct output is generated by the model for a specific theorem, improvement is defined to be zero. We define \textit{nonempty improvement} as the improvement restricted to theorems for which some output has nonzero improvement. Intuitively, improvement is the expected improvement in metric score from the input to output, accounting for errors in the generation. The nonempty improvement score is the expected improvement in metric score, given that there are no errors in the generation.

Additionally, the \textit{accuracy} is the percentage of theorems in the dataset which the model was able to generate a correct output for. The \textit{improved accuracy} is the percentage of theorems in the dataset which the model was able to generate a correct output for, as well as improve the metric to be nonzero.

\subsection{Ablation Studies}
\label{sec:improver_ablation}

\subsubsection{Setup}
\label{sssec:ablations}

When performing our ablation studies, we used a fixed dataset (MIL) and metric (length) and varied the parameters of all the features to find the optimal combination. However, as there are over $8640$ possible combinations, rather than test all combinations, we evaluate using a factorial testing method.

\paragraph{Testing Groups.} We define the following testing groups with the specified parameter combinations:

\textit{GPT-4o-mini/GPT-4o:} This varies the GPT-4o model, outputting a \texttt{string} with no other features.

\textit{Output and CoS:} We evaluate the effects of different output formatting styles (\texttt{string}, \texttt{string list}, \texttt{string tree}) and CoS (True, False), with the model fixed as GPT-4o, with no other features enabled.

\textit{Example Retrieval:} We evaluate the effects of increasing the number of examples provided (multi-shot prompting) in the range of $0,3,5,7,$ and $10$, with the model fixed as GPT-4o, CoS and output formatting fixed as the best combination from the previous test, and no other features enabled.

\textit{Sampling Method:} Here, we evaluate the effects of best-of-$n$ and refinement for a fixed $n=5$. Additionally we test on the refinement cases if forwarding the most recent iteration result, or all previous iteration results is the best, and if we should keep the best out of the iterations, or the most recent. The model is fixed as GPT-4o, CoS, output formatting, and examples are fixed as the best combination from the previous test, and no other features enabled.

\textit{$n$ and Model:} Here, we evaluate the effects of larger $n$ values and different models. We test $n=3,5,7,10,15$ on GPT-4o and GPT-4o-mini, as well as $n=20$ for GPT-4o-mini (as it is of a far lower token cost). CoS, output formatting, examples, and sampling method are fixed as the best combination from the previous test, and no other features enabled.

\textit{Combos and RAG:} We evaluate combination methods \verb|refinement(best_of_m',m)| and \verb|best_of_m'(refinement(m))|, for $m\neq m'$ with $mm'$ equal to the optimal value $m$ from the previous test. We also test the effect of enabling document retrieval. Model, CoS, output formatting, examples, $n$, and sampling method are fixed as the best combination from the previous test.

\paragraph{Selection Criterion.} For each testing group, we select the best parameter combination --- which is then held as constant for the testing of all future testing groups --- based on the combination that has the maximal improvement score. This improvement score represents the expected improvement in metric score, accounting for possible errors in the generation; selecting the parameter combination with the highest such score allows for rewarding both generation accuracy and large improvements in the metric score.

Comparing this with the other three performance metrics, accuracy is not preferred as a selection heuristic, as by simply returning the initial input, we can get $100\%$ accuracy. Improved accuracy accounts for this by only counting theorems that have some positive improvement in metric score in the calculation, but this does not reward larger improvements to metric score any differently than smaller ones. Conversely, nonempty improvement ignores incorrect generations, so it is also not preferable for selection. The improvement score accounts for all this, rewarding correct generations and discouraging incorrect ones, and placing a higher weight to larger improvements in metric score.

\subsubsection{Ablation Results}
\label{ssec:ablation_results}

We perform ablation studies using a subset of the MIL dataset as discussed in \S\ref{sssec:ablations}. The results of this factorial study are aggregated in Table~\ref{tab:abl}.

\begin{table}[h]
\caption{Ablation results. Each cell in the ablation tests shows \texttt{best} / \texttt{worst}, which are the \texttt{best} and \texttt{worst} parameter combinations in the test group.}
\label{tab:abl}
\centering
\small
\begin{tabular}{l|cccc}
\toprule
& Improvement & Nonempty Improve. & Accuracy & Improved Acc.\\
\midrule
GPT-4o-mini & 0 & 0 & 3.62\% & 0\%\\
GPT-4o & 7.03 & 19.67 & 35.77\% & 15.33\%\\
+ Output and CoS & 8.04 / 6.31 & 12.38 / 14.17 & 64.96\% / 44.53\% & 21.17\% / 16.06\% \\
+ Example Retrieval & 9.34 / 5.67 & 14.7 / 8.44 & 63.5\% / 67.15\% & 21.9\% / 16.79\% \\
+ Sampling Method & 15.35 / 9.34 & 18.44 / 14.7 & 83.21\% / 63.5\% & 36.5\% / 21.9\% \\
+ $n$ and Model & 23.51 / 3.65 & 26.28 / 4.63 & 89.47\% / 78.95\% & 45.61\% / 8.77\% \\
+ Combos and RAG & 34.88 / 28.25 & 57.56 / 33.48 & 60.61\% / 84.38\% & 54.55\% / 53.12\% \\
\midrule
\textbf{ImProver} & \textbf{34.88} & \textbf{57.56} & \textbf{100\%} & \textbf{54.55\%}\\
\bottomrule
\end{tabular}
\end{table}

We measure the baseline results from the GPT-4o and GPT-4o-mini models, noting that GPT-4o is the better-scoring model (with respect to the improvement score). Thus, fixing this model, we vary the output formatting type and if CoS is enabled, and determine that outputting \texttt{string list} with CoS enabled maximizes the improvement score. Fixing these parameters, we now vary the number of examples retrieved, noting that prompting with $10$ examples maximizes the improvement score. Fixing this parameter, we vary the sampling methods (excluding compound methods and fixing $n=5$) and observe that best-of-$n$ is the best parameter combination. Now, as GPT-4o-mini is significantly less computationally expensive than its GPT-4o counterpart, we test both models with the sample method fixed to best-of-$n$, and vary $n=1,3,5,7,10,15$, and for GPT-4o-mini, also $n=20$. We conclude that GPT-4o with $n=15$ is the most effective. Fixing these parameters, we consider all mixed compound sampling methods with and without document retrieval enabled, concluding that a $5$-step refinement with best-of-$3$ on each iteration, with RAG enabled, is the optimal combination.

Thus, the optimal parameter combination is GPT-4o outputting as a \texttt{string list} with CoS, RAG, $10$ examples, $5$-step refinement with each iteration being a best-of-$3$ evaluation. Changing any one of these parameters leads to a reduction in performance.

We now present the detailed ablation tables.

\begin{table}[ht]
\caption{Output and Chain-of-States Ablations.}
\label{tab:abl-cos}
\centering
\small
\begin{tabular}{ll|cccc}
\toprule
Output Format & CoS & Improvement & Nonempty Improve. & Accuracy & Improved Acc.\\
\midrule
\texttt{string} & True & 7.53 & 16.12 & 46.72\% & 16.79\% \\
\texttt{string} & False & 7.03 & 19.67 & 35.77\% & 15.33\% \\
\texttt{string list} & \textbf{True} & \textbf{8.04} & \textbf{12.38} & \textbf{64.96\%} & \textbf{21.17\%} \\
\texttt{string list} & False & 7.04 & 13.58 & 51.82\% & 18.98\% \\
\texttt{string tree} & True & 7.62 & 15.34 & 49.64\% & 18.25\% \\
\texttt{string tree} & False & 6.31 & 14.17 & 44.53\% & 16.06\% \\
\bottomrule
\end{tabular}
\end{table}

By Table~\ref{tab:abl-cos}, the optimal combination in this testing group is a \texttt{string list} output format with CoS enabled.

\begin{table}[h]
\caption{Example Retrieval Ablations.}
\label{tab:abl-ex}
\centering
\small
\begin{tabular}{l|cccc}
\toprule
Examples & Improvement & Nonempty Improve. & Accuracy & Improved Acc.\\
\midrule
0 & 5.67 & 8.44 & 67.15\% & 16.79\% \\
3 & 8.49 & 13.68 & 62.04\% & 19.71\% \\
5 & 8.38 & 12.9 & 64.96\% & 21.17\% \\
7 & 7.56 & 12.04 & 62.77\% & 19.71\% \\
\textbf{10} & \textbf{9.34} & \textbf{14.7} & \textbf{63.5\%} & \textbf{21.9\%} \\
\bottomrule
\end{tabular}
\end{table}

By Table~\ref{tab:abl-ex}, the optimal combination in this testing group is $10$ examples.

\begin{table}[h]
\caption{Sampling Method Ablations.}
\label{tab:abl-samp}
\centering
\small
\begin{tabular}{lll|cccc}
\toprule
Method & Forward & Keep Best & Improvement & Nonempty Improve. & Accuracy & Improved Acc.\\
\midrule
None & N/A & N/A & 9.34 & 14.7 & 63.5\% & 21.9\% \\
refinement & 1 & False & 14.76 & 30.63 & 48.18\% & 30.66\% \\
refinement & 5 & False & 12.5 & 20.88 & 59.85\% & 30.66\% \\
refinement & 1 & True & 14.95 & 14.95 & 100.0\% & 30.66\% \\
refinement & 5 & True & 13.15 & 13.15 & 100.0\% & 29.93\% \\
\textbf{best-of-$n$} & N/A & N/A & \textbf{15.35} & \textbf{18.44} & \textbf{83.21\%} & \textbf{36.5\%} \\
\bottomrule
\end{tabular}
\end{table}

Note that the forward and keep-best values are parameters for refinement controlling how many previous iterations to forward, and whether to keep the most recent or the best iteration in subsequent refinement steps. By Table~\ref{tab:abl-samp}, the optimal combination in this testing group is best-of-$n$.

\begin{table}[h]
\caption{Model and $n$ Ablations.}
\label{tab:abl-n}
\centering
\small
\begin{tabular}{ll|cccc}
\toprule
Model & $n$ & Improvement & Nonempty Improve. & Accuracy & Improved Acc.\\
\midrule
gpt-4o & 3 & 19.66 & 24.36 & 80.7\% & 38.6\% \\
gpt-4o & 5 & 20.12 & 24.97 & 80.56\% & 36.11\% \\
gpt-4o & 7 & 22.44 & 27.21 & 82.46\% & 42.11\% \\
gpt-4o & 10 & 21.73 & 25.28 & 85.96\% & 40.35\% \\
\textbf{gpt-4o} & \textbf{15} & \textbf{23.51} & \textbf{26.28} & \textbf{89.47\%} & \textbf{45.61\%} \\
gpt-4o-mini & 3 & 3.65 & 4.63 & 78.95\% & 8.77\% \\
gpt-4o-mini & 5 & 5.12 & 6.21 & 82.46\% & 10.53\% \\
gpt-4o-mini & 7 & 3.65 & 4.34 & 84.21\% & 8.77\% \\
gpt-4o-mini & 10 & 4.99 & 5.69 & 87.72\% & 12.28\% \\
gpt-4o-mini & 15 & 4.35 & 5.06 & 85.96\% & 12.28\% \\
gpt-4o-mini & 20 & 4.87 & 5.56 & 87.72\% & 14.04\% \\
\bottomrule
\end{tabular}
\end{table}

By Table~\ref{tab:abl-n}, the optimal combination in this testing group is GPT-4o with $n=15$.

\begin{table}[h]
\caption{RAG and Combination Sampling Method Ablations.}
\label{tab:abl-rag}
\centering
\small
\begin{tabular}{llll|cccc}
\toprule
Combination & $m$ & $m'$ & RAG & Improvement & Nonempty Improve. & Accuracy & Improved Acc.\\
\midrule
best-of-$n$(refinement) & 3 & 5 & True & 33.78 & 33.78 & 100.0\% & 50.0\% \\
best-of-$n$(refinement) & 3 & 5 & False & 31.23 & 31.23 & 100.0\% & 46.88\% \\
best-of-$n$(refinement) & 5 & 3 & True & 31.85 & 31.85 & 100.0\% & 50.0\% \\
best-of-$n$(refinement) & 5 & 3 & False & 31.35 & 31.35 & 100.0\% & 50.0\% \\
refinement(best-of-$n$) & 3 & 5 & True & 32.66 & 51.32 & 63.64\% & 48.48\% \\
refinement(best-of-$n$) & 3 & 5 & False & 32.88 & 50.1 & 65.62\% & 53.12\% \\
\textbf{refinement(best-of-$n$)} & \textbf{5} & \textbf{3} & \textbf{True} & \textbf{34.88} & \textbf{57.56} & \textbf{60.61\%} & \textbf{54.55\%} \\
refinement(best-of-$n$) & 5 & 3 & False & 29.54 & 49.75 & 59.38\% & 43.75\% \\
best-of-$n$ & N/A & 15 & True & 29.64 & 32.71 & 90.62\% & 56.25\% \\
best-of-$n$ & N/A & 15 & False & 28.25 & 33.48 & 84.38\% & 53.12\% \\
\bottomrule
\end{tabular}
\end{table}

By Table~\ref{tab:abl-rag}, the optimal combination in this testing group is a $5$-step refinement with each iteration being a best-of-$3$ call, with RAG enabled.

\paragraph{Declarativity and Chain-of-States Ablation.} We additionally examine the effects of disabling CoS on declarativity optimization tasks, as we speculate that CoS has a high impact on the performance of declarativity optimization tasks, since the proof states that are embedded due to CoS seem to be a critical aspect to generating the explicit declarations that the declarative metric measures.

\begin{table}[t]
\caption{CoS Declarativity Ablation results.}
\label{tab:read-abl}
\centering
\small
\begin{tabular}{l|cccc}
\toprule
& Improvement & Nonempty Improve. & Accuracy & Improved Acc.\\
\midrule
GPT-4o & 4.97 & 15.89 & 37.5\% & 12.5\% \\
ImProver, CoS Disabled & 9.23 & 24.61 & 100.0\% & 28.12\% \\
\textbf{ImProver} & \textbf{16.69} & \textbf{31.42} & \textbf{100.0\%} & \textbf{46.88\%} \\
\bottomrule
\end{tabular}
\end{table}

We confirm this result by considering Table~\ref{tab:read-abl} and observe that enabling CoS nearly doubles the improvement score, and significantly improves the nonempty improvement score, suggesting that CoS has a large impact on optimizing for the declarative metric, as conjectured. However, we also note a significant increase in improved accuracy, which suggests that embedding the chain of states also improves the ability of the model to generate nontrivial correct outputs, implying that the symbolic information contained in the states is critical to effectively making a proof more declarative.

\paragraph{Syntax Guidance Ablation.} We examine the effects of syntax guidance on ImProver's performance. To test this, we consider a subset of MIL, and optimize for length with and without error message forwarding.

\begin{table}[t]
\caption{Syntax Guidance Ablation results.}
\label{tab:syn-abl}
\centering
\small
\begin{tabular}{l|cccc}
\toprule
& Improvement & Nonempty Improve. & Accuracy & Improved Acc.\\
\midrule
GPT-4o & 11.00 & 25.94 & 42.42\% & 21.21\% \\
ImProver, No Syntax Guidance & 23.42 & 49.97 & 100.0\% & 46.88\% \\
\textbf{ImProver} & \textbf{28.94} & \textbf{48.74} & \textbf{100.0\%} & \textbf{59.38\%} \\
\bottomrule
\end{tabular}
\end{table}

Considering the results of this ablation in Table~\ref{tab:syn-abl}, we observe that without syntax guidance and error forwarding, the ability of the model to improve the metric score is approximately unchanged, but there is a significant $13\%$ spike in improved accuracy. This signifies that the syntax guidance improves the model's ability to generate correct results --- as is expected --- but does not improve the model's ability to optimize proofs assuming correct generations. This ensures that the large improvement in performance compared to GPT-4o is not solely due to simple syntax guidance, but is more so caused by improvements like CoS, example retrieval, document retrieval, and so on.

\subsection{Main Results}
\label{sec:improver_results}

\begin{table}[t]
\caption{Average proof optimization results aggregated across all datasets.}
\label{tab:main-avg}
\centering
\small
\begin{tabular}{ll|cccc}
\toprule
Metric & Model & Improvement & Nonempty Improvement & Accuracy & Improved Acc.\\
\midrule
\multirow{2}{*}{\textbf{Length}} & GPT-4o & 3.7 & 15.15 & 26.36\% & 8.31\% \\
 & \textbf{ImProver} & \textbf{20.96} & \textbf{55.29} & \textbf{100.0\%} & \textbf{35.44\%} \\
\midrule
\multirow{2}{*}{\textbf{Declarativity}} & GPT-4o & 2.21 & 8.02 & 18.75\% & 6.13\% \\
 & \textbf{ImProver} & \textbf{9.34} & \textbf{30.53} & \textbf{100.0\%} & \textbf{24.56\%} \\
\midrule
\multirow{2}{*}{\textbf{Mixed}} & GPT-4o & 3.51 & 23.90 & 14.70\% & 5.11\% \\
 & \textbf{ImProver} & \textbf{27.31} & \textbf{39.24} & \textbf{100.0\%} & \textbf{30.55\%} \\
\bottomrule
\end{tabular}
\end{table}

\begin{table}[t]
\caption{MIL proof optimization results.}
\label{tab:main-mil}
\centering
\small
\begin{tabular}{ll|cccc}
\toprule
Metric & Model & Improvement & Nonempty Improvement & Accuracy & Improved Acc.\\
\midrule
\multirow{2}{*}{\textbf{Length}} & GPT-4o & 6.25 & 18.58 & 37.5\% & 14.42\% \\
 & \textbf{ImProver} & \textbf{30.54} & \textbf{56.56} & \textbf{100.0\%} & \textbf{50.0\%} \\
\midrule
\multirow{2}{*}{\textbf{Declarativity}} & GPT-4o & 4.18 & 14.48 & 28.85\% & 11.54\% \\
 & \textbf{ImProver} & \textbf{13.45} & \textbf{30.97} & \textbf{100.0\%} & \textbf{34.21\%} \\
\midrule
\multirow{2}{*}{\textbf{Mixed}} & GPT-4o & 3.70 & 27.38 & 13.51\% & 0.0\% \\
 & \textbf{ImProver} & \textbf{43.55} & \textbf{44.76} & \textbf{100.0\%} & \textbf{45.94\%} \\
\bottomrule
\end{tabular}
\end{table}

\begin{table}[t]
\caption{Compfiles proof optimization results.}
\label{tab:main-compfiles}
\centering
\small
\begin{tabular}{ll|cccc}
\toprule
Metric & Model & Improvement & Nonempty Improvement & Accuracy & Improved Acc.\\
\midrule
\multirow{2}{*}{\textbf{Length}} & GPT-4o & 2.75 & 30.7 & 11.54\% & 5.13\% \\
 & \textbf{ImProver} & \textbf{18.86} & \textbf{54.48} & \textbf{100.0\%} & \textbf{34.62\%} \\
\midrule
\multirow{2}{*}{\textbf{Declarativity}} & GPT-4o & 0.39 & 3.38 & 14.1\% & 1.28\% \\
 & \textbf{ImProver} & \textbf{5.74} & \textbf{24.89} & \textbf{100.0\%} & \textbf{19.23\%} \\
\midrule
\multirow{2}{*}{\textbf{Mixed}} & GPT-4o & 3.96 & 3.96 & 26.9\% & 20.0\% \\
 & \textbf{ImProver} & \textbf{20.60} & \textbf{76.53} & \textbf{100.0\%} & \textbf{23.07\%} \\
\bottomrule
\end{tabular}
\end{table}

\begin{table}[h]
\caption{Mathlib proof optimization results.}
\label{tab:main-lib}
\centering
\small
\begin{tabular}{ll|cccc}
\toprule
Metric & Model & Improvement & Nonempty Improvement & Accuracy & Improved Acc.\\
\midrule
\multirow{2}{*}{\textbf{Length}} & GPT-4o & 0.0 & 0.0 & 16.67\% & 0.0\% \\
 & \textbf{ImProver} & \textbf{6.19} & \textbf{53.65} & \textbf{100.0\%} & \textbf{11.54\%} \\
\midrule
\multirow{2}{*}{\textbf{Declarativity}} & GPT-4o & 0.0 & 0.0 & 4.65\% & 0.0\% \\
 & \textbf{ImProver} & \textbf{4.63} & \textbf{33.19} & \textbf{100.0\%} & \textbf{11.63\%} \\
\midrule
\multirow{2}{*}{\textbf{Mixed}} & GPT-4o & 2.92 & \textbf{30.14} & 9.30\% & 4.65\% \\
 & \textbf{ImProver} & \textbf{4.16} & 7.45 & \textbf{100.0\%} & \textbf{9.30\%} \\
\bottomrule
\end{tabular}
\end{table}

\textbf{ImProver is capable of optimizing proofs in all settings.} From Table~\ref{tab:main-mil}, Table~\ref{tab:main-compfiles}, and Table~\ref{tab:main-lib}, we can see that ImProver is capable of optimizing proofs on all datasets for both the length and declarative metrics, as well as on the mixed metric. Furthermore, Table~\ref{tab:main-avg} shows that across all metrics, ImProver significantly outperforms GPT-4o on proof optimization tasks on every experimental measure --- aggregated from all datasets. Additionally, from Table~\ref{tab:main-mil}, Table~\ref{tab:main-compfiles}, and Table~\ref{tab:main-lib}, we can see that ImProver outperforms GPT-4o on each dataset as well.

\textbf{Length optimization.} First focusing on the length metric, we see that ImProver outperforms GPT-4o with respect to the improvement score by $566\%$ (aggregated over all datasets). Additionally, we are guaranteed that ImProver produces a correct output, although that output may just be the same as the input. However, $35.44\%$ of the time, it generates a correct output that is not the same length as the input, and in that case, we expect an average of a $55.29\%$ reduction in length. Comparing this with GPT-4o, we conclude that not only can ImProver optimize at a higher level on arbitrary theorems, but its ability to generate nontrivial correct outputs is far greater in comparison to GPT-4o.

\textbf{Declarativity optimization.} Declarativity optimization is similar, with ImProver outperforming GPT-4o by $423\%$. Moreover, the accuracy, improved accuracy, and nonempty improvement disparities for declarativity parallel those of the length tests. However, it should be noted that for both GPT-4o and ImProver, the accuracy and improved accuracy scores were markedly smaller for declarativity than length optimization. This suggests that for both models, it was generally more ``difficult'' to generate a correct output, and moreover, generate a correct output with a better metric score than the input, for declarativity optimization than length optimization. In other words, optimizing for declarativity is more difficult for the underlying generator than optimizing for length. However, we speculate with higher-quality prompts and metrics, this disparity can be minimized. Regardless, we note that different metrics can be less likely to be correctly optimized, and that model performance is correlated with the metric it seeks to optimize, both for GPT-4o and ImProver.

\textbf{Mixed optimization.} For the mixed optimization of both length and declarativity, we observe a $778\%$ outperformance by ImProver in the aggregate, with similar increases across all experimental measures. It is notable that the combined improvements and accuracies mirror the trends from the individual length and declarativity metrics, suggesting that ImProver is able to scale its optimizations to align with the increased complexity of the metric. Moreover, this empirically shows that for more complex and arbitrarily-defined metrics, ImProver maintains its ability to generate nontrivial optimizations, given sufficiently high-quality examples, prompts, and scoring functions.

\textbf{Optimization varies based on dataset difficulty.} Additionally noting Table~\ref{tab:main-mil}, Table~\ref{tab:main-compfiles}, and Table~\ref{tab:main-lib}, we observe that the improvement score for all metrics for both GPT-4o and ImProver is highest for the MIL dataset, lower for Compfiles, and the lowest on the Mathlib theorems. This suggests that the expected improvement in metric score decreases with higher difficulty, with undergraduate-level theorems having a significantly higher expected improvement than research-level theorems. However, it should be noted that for all metrics, the nonempty improvement of ImProver stayed somewhat consistent, whereas for GPT-4o, it followed the aforementioned trend of decreasing with difficulty. Similarly, the accuracy and improved accuracy scores for both metrics and models decreased with higher difficulty datasets (disregarding ImProver's accuracy scores, as they are ensured to be $100\%$). This suggests that although the base GPT-4o generator is less likely to generate a correct output for higher difficulty datasets, the improvements that ImProver makes to the base generator allow it to maintain its improvement in the metric score whenever a correct output is generated. As such, we can speculate that the bottleneck in the improvement score is not the model's ability to optimize the proof for a metric, but rather its ability to generate a new correct proof at all. As such, we conjecture that with more capable generator models, the accuracy --- and thus, the improvement score --- in optimization tasks will continue to increase, until the improvement scores match the nonempty improvement.

Overall, we conclude that although the performance of both ImProver and GPT-4o decreases on length and declarativity optimization on more difficult datasets, ImProver significantly outperforms GPT-4o on all datasets for length, declarativity, and mixed optimization.

\subsection{Neural Theorem Proving Evaluation}
\label{ssec:ntp}

We evaluate ImProver's neural theorem proving (NTP) performance using the completion metric on a subset of MIL with empty input proofs. Table~\ref{tab:comp_ntp} shows the accuracy on the dataset split by topic for both ImProver and GPT-4o.

\begin{table}[ht]
\caption{Proof generation results. Each cell shows percent accuracy.}
\label{tab:comp_ntp}
\centering
\small
\begin{tabular}{l|ccc|c}
\toprule
& MIL-C04 Pass@15 & MIL-C08 Pass@15 & MIL Pass@15 & MiniF2F-test Pass@8 \\
\midrule
GPT-4o & 18.18\% & 25\% & 21.73\% & 9.02\% \\
ImProver & \textbf{45.45\%} & \textbf{33.33\%} & \textbf{39.13\%} & 16.39\% \\
Lean Expert Iteration & --- & --- & --- & \textbf{34.5\%} \\
\bottomrule
\end{tabular}
\end{table}

ImProver substantially outperforms GPT-4o across all datasets. Thus, proof optimization systems do indeed generalize NTP systems: optimizing for the completion metric (correctness) is equivalent to neural theorem proving.

Additionally, we note that specialized neural theorem proving models such as GPT-f \cite{zheng2022miniff} and Lean Expert Iteration \cite{DBLP:journals/corr/abs-2202-01344} outperform ImProver on MiniF2F, but as they are specially trained on a large corpus of Lean code (whereas ImProver is an extension of GPT-4o), as well as designed for theorem proving rather than proof optimization, it is worthwhile for future works to consider fine-tuning or specializing ImProver's methodology specifically for theorem proving. The purpose of this experiment is to empirically prove that proof optimization systems like ImProver do indeed generalize the problem of theorem proving as claimed in \S\ref{ssec:metrics_31}.

\subsection{Discussion}
\label{sec:improver_discussion}

ImProver demonstrates that automated proof optimization is achievable for undergraduate, competition, and research-level mathematics. The system's core strengths are:

\begin{enumerate}
    \item \textbf{Correctness guarantee}: By always returning at least the original proof as a fallback, ImProver achieves $100\%$ accuracy.
    \item \textbf{Chain-of-States}: The symbolic intermediate state information is the single most impactful improvement, especially for structural metrics.
    \item \textbf{Retrieval}: Both example and document retrieval provide critical context, particularly for finding relevant Mathlib lemmas.
    \item \textbf{Error correction}: Iterative refinement substantially improves the quality of generated proofs.
\end{enumerate}

However, ImProver is limited by its high computational cost, which is exacerbated by its reliance on proprietary LLMs (GPT-4o). In future work, we will apply supervised fine-tuning and reinforcement learning on a smaller model to match performance locally. ImProver demonstrates its ability to generate substantially shorter and more declarative proofs while maintaining correctness. As such, we believe that ImProver sets the stage for further work on proof optimization to advance the study and use of AI in mathematics.
